\subsection{The completed irreducible classification and sign theorem}

\begin{statement}{P4.58}{Proved}
Let $f:Y\to B$ be the fibration of Part~4.2, with
$b_1(Y)=b_2(Y)=0$, connected fibres and general fibre a
complex two-torus. Suppose every reduced fibre support
is irreducible or has normal crossings. Then
\[
 B=\mathbb P^1,\qquad c_3(Y)\in\{0,2\},
                       \qquad b_3(Y)=2-c_3(Y).
\]
Every reduced support has nonnegative Euler number.
An irreducible singular support whose normalization has
a rational resolution has Euler number two. Every other
irreducible support has Euler number zero. Such a support
is either smooth with a finite unramified torus cover,
or singular with smooth minimal elliptic ruled
normalization, or singular with smooth minimal Hopf
normalization. In particular the class-VII normalization
alternative is realized only by a smooth minimal Hopf
surface.

For a normal-crossing support the Euler number equals
its number of triple points, and is even. Its zero case
is exactly P4.8: in the reducible case the normalization
is a disjoint union of smooth minimal Hopf or non-K\"ahler
elliptic surfaces, with smooth elliptic conductor curves
and finite unramified branch maps. The equality
$c_3(Y)=0$ holds exactly when every support has one of
the corresponding zero-Euler forms.
\end{statement}
\begin{proof}
P4.26 proves the sign for every irreducible support,
including singular normalizations, and P4.6--P4.8 give
the full normal-crossing calculation. P4.35 proves the
exact value two for the rationally resolved singular
case. P4.36--P4.37 show that every class-VII case has
smooth minimal Hopf normalization and Euler number
zero. P4.32 gives the complete irreducible zero-Euler
classification and its converses.

We finish the positive-genus ruled alternative without
a transverse-type assumption on its vertical conductor
branches. P4R.13 and the preceding reductions show that
a positive-Euler ruled support would have a doubled
horizontal pole section and an image conductor with two
rational normalized components and one graph cycle.
P4.54 excludes contraction of that section, using the
actual weighted normalization and blowdown maps.
P4R.18 reduces its generic image to $D_4$ with degree
two or three. P4.55 makes that image smooth and excludes
degree three, leaving the degree-two section and one
rational vertical third sheet. P4.56 forces at least
three generic branches on the other image component.

The graph alternatives are exhaustive by P4.55. P4.57
excludes two distinct junctions. P4R.20 excludes a sole
three-branch image point, with the exact conductor-path
and endpoint normalization calculation. Thus no
positive-Euler ruled support remains. This proves the
asserted value zero for every nonrational irreducible
support, and P4.32 supplies its actual normalization.

Finally Euler additivity gives $c_3(Y)$ as the sum of
these nonnegative support Euler numbers. The base,
Poincar\'e-duality calculation $e(Y)=2-b_3(Y)$, and
parity are proved in P4.26. They yield the displayed
global values. A sum of nonnegative integers vanishes
exactly when every summand vanishes, giving the equality
classification through P4.32 and P4.8.

The hypothesis checks are the explicit constructions
P4.10 and P4.11. Removing $b_2=0$ permits the proved
Euler value $24$. Removing the torus-fibre hypothesis
permits the Hopf-fibre example, which fails the stated
normalization equality classification. The numerical
normal-crossing sign itself follows from the stronger
calculation in P4.7 and does not require that hypothesis.
\end{proof}
